\section{Reproducible notebooks and figures}

This paper is accompanied by a set of computational notebooks that generate all figures.
These notebooks are expository: they illustrate the structure of the proof pipeline but do not replace the formal estimates.

\subsection{Notebook structure}

\begin{itemize}
\item \texttt{01\_decay\_weight\_visualization.ipynb}  
Illustrates two-time decay and Carleman-style exponential weighting.

\item \texttt{02\_hardy\_gate\_demo.ipynb}  
Shows how concentration incurs variation cost via Hardy-type constraints.

\item \texttt{03\_contradiction\_pipeline\_pde\_aligned.ipynb}  
Combines weighted norms, variation cost, and curvature proxies to demonstrate the contradiction structure.
\end{itemize}

\subsection{Figure provenance}

Each figure in the paper is generated directly from the notebooks:

\begin{center}
\begin{tabular}{ll}
Figure & Source notebook \\
Step 1 & Notebook 01 \\
Step 2 & Notebook 01 \\
Step 3 & Notebook 02 \\
Step 4 & Notebook 03 \\
Step 5 & Notebook 03 \\
\end{tabular}
\end{center}

\subsection{Interpretation vs proof}

The quantities computed in the notebooks (weighted norms, gradient costs, curvature proxies) are finite-grid analogues of terms appearing in the formal analysis.

They are used to:

\begin{itemize}
\item visualize constraint interactions,
\item illustrate incompatibility for non-zero candidates,
\item support the logical structure of the contradiction argument.
\end{itemize}

They are not substitutes for the analytical estimates in \cite{EKPV2008, TaoPDE}.

\subsection{Access}

The notebooks and figures are available in the repository:

\begin{center}
\texttt{https://github.com/thinkthoughts/unique-continuation-constraint-lab}
\end{center}

\subsection{Summary}

The appendix connects:

\[
\text{analysis} \leftrightarrow \text{computation} \leftrightarrow \text{visualization}.
\]

This alignment is intended to improve clarity without altering the underlying mathematics.
